
Contamination detection for foundation model evaluation benchmarks is complicated by inconsistent signals across detection methods applied to the same benchmarks. This paper describes an attempt to build a geometry-based routing system — a three-zone phase diagram — that classifies benchmark items by corpus-side overlap signals (max 13-gram count and SBERT cosine similarity to nearest corpus neighbor) and routes each item to the detector family best suited to its contamination geometry. The experiment evaluates 25,403 items from MMLU (14,042), HellaSwag (10,042), and GSM8K (1,319) against three pretraining corpora (The Pile, C4, FineWeb) using five detector families. The primary result is not the routing accuracy the experiment was designed to measure but rather the identification of a failure mode called stratum collapse: random corpus streaming produces near-zero SBERT cosine similarity distributions at 50,000-document scale, collapsing all items into a single lexical stratum and making routing structurally impossible. Top-k nearest-neighbor retrieval is identified as a necessary design prerequisite for semantic stratum formation. Two secondary empirical findings are confirmed with cross-corpus reproducibility: MMLU achieves n-gram recall = 1.0 against all three corpora at 50K-document sampling, and GSM8K achieves recall = 0.0. Four implementation bugs in the detector and evaluator code are identified and documented. The primary routing hypothesis (P1) is inconclusive; gate metrics for n-gram lexical recall (0.556 vs. target ≥ 0.80), Min-K\%++ F1 variance (0.000 vs. target ≥ 0.15), and indeterminacy rate (0.000 vs. target [0.10, 0.50]) are not met.

